% Список пользовательских команд 

%%%%%%%%%%%% \image %%%%%%%%%%%%%%%%%%

% \image {Имя изображения.расширение}{Подпись к рисунку}{Скейл Изображения}

\newcommand{\image}[3]{
\begin{figure}[H]
	\centering
	\includegraphics[width=#3\textwidth]{#1}
	\caption{#2}
\end{figure}
}


%%%%%%%%%%%%%%%%%%%%%%%%%%%%%%%%%%%%%%



%%%%%%%%%% \codefromfile  %%%%%%%%%%%%%%%%%%

% \codefromfile {Имя файла}

\newcommand{\codefromfile}[2]{
\begin{code}
	\inputminted[breaklines=true, framesep=10pt, fontsize=\footnotesize, firstline=1,]{#2}{Listings/#1}
\end{code}
}


%%%%%%%%%%%%%%%%%%%%%%%%%%%%%%%%%%%%%%



%%%%%%%%%%%% Окружение algorithm %%%%%%%%%%%%%%%%%

% Используется в content.tex как:
%   \begin{algorithm}{Название}\hfill\\
%       строки псевдокода со \\ в конце
%   \end{algorithm}

\newenvironment{algorithm}[1]
  {%
    \par\medskip\noindent\textbf{Алгоритм (#1).}%
    \par\nopagebreak\smallskip
    \begin{flushleft}\setlength{\parindent}{0pt}%
  }
  {%
    \end{flushleft}\par\medskip
  }

%%%%%%%%%%%%%%%%%%%%%%%%%%%%%%%%%%%%%%